%%
%% Author: shoupu_wan
%% 7/14/18
%%

\chapter{Organize your Conditional Statements}\label{ch:organizeYourIfElse}

Conditional is a control switching mechanism to alter code execution path,
conditioned by boolean expressions.
As a control flow, conditional is one of the fundamental building blocks
of almost all programming languages.
Undoubtedly, many references---textbooks, research articles, blog posts---have already covered this
comprehensive discussions
\footnote{\url{https://en.wikipedia.org/wiki/Conditional_(computer_programming)}}.
It may seem strange that I am spending a chapter on such a `well-known' topic here and now.
But based on my years of working experience in several top-tier tech companies,
haphazard suboptimal solutions with convoluted and lengthy logic make their way to
production codebase and becomes the
core for millions of businesses.
The crooked design, the obscure the logic later pose as road-block and hazard on the way of
code refactor and reuse.
This shows that just because the topic is basic and `well-known' does not mean that folks can
or will do it right.
%But it is seldom the case that the first version of implementation of any non-trivial logic, if
%correct at all, is optimal.
``As long as it works, why does one bother to simplify the logic?'', some may argue.
My answer is ``Suboptimal, obscure logic fails the primary purpose of coding---communication.''
Worse, it often results in redundant/duplicate code which lays outrageous burden on code
maintenance.
I'd like to quote Steve McConnell
\begin{quote}
    ``..the primary audience [of computer programs] is made up of humans,
    not computers\cite{Mcconnell.2004}.''
\end{quote}
Martin Fowler put it even more straightforward
\begin{quote}
    ``Any fool can write code that a computer can understand.
    Good programmers write code that humans can understand\cite{Fowler1999}.''
\end{quote}

Twisted code and intertwined logic

I don't blame the author of twisted code and intertwined logic, which
are often the victorious output of a fierce battle fought by gallant developers.
They has all my respect.
But I criticize anyone who intends half-cooked work for final product in defiance of the
production standard.
Unequivocally I advocate that developers finish the other half of the work
---polish the logic, simplify the conditional statements, craft production-quality solutions.
Developers finish walking the last mile, solidify your otherwise hard-won victory.

With simple techniques and due diligence, producing beautiful, succinct solution
with initially unrefined code is possible.
The rest of this chapter is dedicated to this.

%Communicate clear, succinct code to colleagues and successors are much more efficient.
The techniques are based on my years of software development experience.
%of how to polish complicated logic involving conditionals into simple, beautiful, and succinct code.
You will learn how to design conditional statements using \code{if}, \code{else
if} (or \code{elif}), \code{else}, \textit{etc.},
to produce effective, succinct logic.
You will learn to minimize occurrences of redundancy condition checking in your code and reduce
duplicate code overall.
Through plenty of examples, I demonstrate how to apply these techniques for logic simplification.
In doing so, I hope programmers would be inspired to strive for beautiful, succinct conditional
statements.
I believe the effort you spend in studying the basics of conditional programming
will soon pay itself off.

\section{Equivalent but simple conditional expressions}\label{sec:logicEquivalences}

A boolean variable or expression can evaluate to two possible values-\code{true} and \code{false}.
Depending on the evaluation result of a boolean value,
a \code{if-else} test has two execution alternatives-the \code{if}-block or
the \code{else}-block;
A composite condition consisting of two boolean values may have up to 4 alternatives.
Three boolean values make have up to 8 alternatives.
And so on.
\begin{figure}[htb]
    \centering
    \includegraphics[width=.7\textwidth]{if-else/nested-conditionals.jpg}
    \caption{\label{fig:conditional-branching}An example of $3$ nested conditional statements.}
\end{figure}

Establishing the equivalence between logic and conditional statements can be subtle.
Simplifying them is even harder, especially when camouflaged in the
jungle of real-world code with long variable names, complex expression, comments, \textit{etc.}
Take a look at ~\autoref{lst:if-return-boolean-real-world}.
At least to me, it is not yet clear how to simplify the logic or even what it does.
\lstinputlisting[caption={Code snippet for demonstration},
label={lst:if-return-boolean-real-world}]
{if-else/equivalent_statements_if_return_real_world.java}
But if we replace the \code{if}-test with \code{<test>} and filter out the
comments, it would look like ~\autoref{lst:if-return-boolean}.
Now it is clear that the code performs a test and return a boolean value accordingly.
You may even realize that it is actually equivalent to ~\autoref{lst:simplified-return-boolean}.
% side-by-side two columns
\newline\begin{minipage}[t]{.45\textwidth}
            \lstinputlisting[caption={Unsimplified test and return boolean code},
            label={lst:if-return-boolean},
            firstline=1,lastline=4]
            {if-else/equivalent_statements_if_return.java}
\end{minipage}\hfill
\noindent\begin{minipage}[t]{.45\textwidth}
             \lstinputlisting[caption={Simplified statement},
             label={lst:simplified-return-boolean},
             firstline=6,lastline=6]
             {if-else/equivalent_statements_if_return.java}
\end{minipage}
But when embedded in real problems, disguised by specific forms, especially when some alternatives
are masked, achieving optimal conditional programming may not be so straightforward.

Many of the techniques employed in this section are of math origin in the domain of Boolean algebra.
I frequently make references to the logical equivalences, and normal forms, truth tables,
\textit{etc.}
Review and refresh your memory on relevant materials as needed.
Let us examine the code in ~\autoref{lst:nestedBranching}\footnote{ this example is adapted
from my own code at earlier years of software development career.}

% side-by-side two columns
\noindent\begin{minipage}[t]{.3\textwidth}
             \lstinputlisting[label={lst:nestedBranching},
             caption={v1}]
             {if-else/composite_condition_tobe_simplified.java}
\end{minipage}\hfill
\noindent\begin{minipage}[t]{.3\textwidth}
             \lstinputlisting[label={lst:simpleBranching},
             caption={v2}]
             {if-else/composite_condition_simplified1.java}
\end{minipage}\hfill
\noindent\begin{minipage}[t]{.3\textwidth}
             \lstinputlisting[label={lst:final-simplified-conditional-statement},
             caption={v3}]
             {if-else/composite_condition_simplified2.java}
\end{minipage}


It is an intermediate version that barely converged on fulfilling all the logical requirements.
As expected, it works but `smells'---it has complicated, nested conditional statements and duplicate
code\cite{Fowler1999}.
The readability is so compromised that one can hardly follow the paths of symbolic execution
without a great deal of mental twists.
\begin{table}[bth]
    \begin{center}
        \caption{Normal forms for conditions of execution of each statement.
        \label{tbl:breakdown-of-conditions}}
        \begin{tabular}{|l|l|}
            \hline
            code & condition to execute \\
            \hline
            statement1 & test1 $\land$ test2 \\
            statement2 & test1 $\land\lnot$ test2 $\lor\lnot$ test1 \\
            statement3 & test1 \\
            \hline
        \end{tabular}
    \end{center}
\end{table}
Nevertheless, this is still a good starting point for further improvement---
just do not settle.
Upon capturing this solution, I immediately headed to eliminate the nested
conditional statements.
Observing that the \code{else} block of \code{if (test1)} is much simpler, I decided to flip
the logic by negating \code{test1} so I can swap the two blocks.
Then I promoted the nested \code{if (test2)} up to the first-level and now they
become sequential conditional statements comprised of \code{if..else if..else} blocks
that are built on top of one another.
The resulting code is as shown in ~\autoref{lst:simpleBranching}.
In fact, we just came across a programming pattern that and that is to be explained later
in ~\autoref{sec:case-stack}.

Comparing with ~\autoref{lst:nestedBranching} the conditional statements are now streamlined and
easier to understand.
But still it has code duplication as both \code{statement2} and \code{statement3} are repeated.
So I did a back-of-envelop calculation---
I obtained the needed conditions in order for \code{statement1}, \code{statement2} and
\code{statement3} to be executed, respectively.

The result is shown in ~\autoref{tbl:breakdown-of-conditions}.
With this table, it is clear that the key to simplify the logic lies in the conditions for the
execution of \code{statement2}.
Let us do some Boolean algebra.
Let $c$ be the condition for \code{statement2} to be executed and
\begin{align*}
    a &= \text{\code{test1}}, \\
    b &= \text{\code{test2}}.
\end{align*}
Then
\begin{align}
    c &= a \land \lnot b \lor \lnot a  \nonumber \\
    &= a \land \lnot b \lor (\lnot a \land \lnot b \lor \lnot a \nonumber) & \text{// absorption law
    }\\
    &= (a \land \lnot b \lor \lnot a \land \lnot b) \lor \lnot a \nonumber & \text{// associativity
    law}\\
    &= \lnot b \lor \lnot a & \text{// distributivity law} \nonumber \\
    &=\lnot (a \land b) & \text{// De Morgan's law} \label{eq:boolean-expression-simplify}
\end{align}

~\autoref{eq:boolean-expression-simplify} shows that the condition of execution for
\code{statement2} is \code{!(test1 && test2)} and this happens to fit as the
\code{else}-block of the conditional statement \code{if (test1 && test2)..}.
With this hard-earned wit, it now becomes straightforward to come up with yet an even simpler
implementation such as ~\autoref{lst:final-simplified-conditional-statement}.
This completes the whole logic simplification process.

Note that there are other ways to achieve the same goal.
One can also use \emph{truth tables} to simplify logic.
Let's see how this works.
\begin{table}[h]
    \begin{center}
        \caption{True value table for \code{test1} and \code{test2}.
        \label{tbl:true-value-table}}
        \begin{tabular}{|c|c|c|}
            \hline
            \code{test1} & \code{test2} & true expr \& execution\\
            \hline
            \code{false} & \code{false} & \code{!test1} $\rightarrow$ \code{statement2} \\
            \code{false} & \code{true} & \code{!test1} $\rightarrow$ \code{statement2} \\
            \code{true} & \code{false} & \code{test1 \&\& !test2}
            $\rightarrow$ \code{statement2} \\
            & & \code{test1=true} $\rightarrow$ \code{statement3} \\
            \code{true} & \code{true} & \code{test1 \&\& test2} $\rightarrow$
            \code{statement1} \\
            & & \code{test1} $\rightarrow$ \code{statement3} \\
            \hline
        \end{tabular}
    \end{center}
\end{table}
I have constructed a truth table with \code{test1} and \code{test2} as the variables
(~\autoref{tbl:true-value-table}).
Since there are two conditional variables, there are a total four rows in the table.
Additionally, we added the third column to shows what \code{true} boolean expression leads
the execution of which statement.
Summarizing the table entries, we arrive at the following conclusions
\begin{enumerate}[label={\roman*:}]
    \item \code{statement1} is executed \emph{iff} both \code{test1} and \code{test2}
    are \code{true};
    \item \code{statement2} is executed \emph{iff} not both \code{test1} and \code{test2} are
    \code{true};
    \item \code{statement3} is executed \emph{iff} \code{test1} is \code{true}, regardless of
    \code{test2}.
\end{enumerate}
With all the mental twists untied, implementation of something like
~\autoref{lst:final-simplified-conditional-statement} naturally follows.

Thus far we have presented two processes for logic simplification.
One of them en routes on the necessary and sufficient conditions for the execution of certain
statements and use that to simplify the conditional expressions.
The other lists all the possible branches in a true table and looks for simple ways to combine
them into as few conditional expressions as possible.
In practice, you are encouraged to try both and form an intuition for each process.
You will find each of them has an edge for certain circumstances.

To verify your understanding of the material, in
~\autoref{lst:exercise-original-conditional-statement} is presented an exercise.
Please work it out with one of the two methods above and my simplified solution is provided in
~\autoref{lst:exercise-simplified-conditional-statement} for your information.

% side-by-side two columns
\noindent\begin{minipage}[t]{.3\textwidth}
             \lstinputlisting[label={lst:exercise-original-conditional-statement},
             caption={Simplification exercise}]
             {if-else/generic_conditioning1.java}
\end{minipage}\hfill
\noindent\begin{minipage}[t]{.3\textwidth}
             \lstinputlisting[label={lst:exercise-simplified-conditional-statement},
             caption={Simplified}]
             {if-else/generic_conditioning2.java}
\end{minipage}

\section{Short-circuit evaluation}
\label{sec:short-circuit-eval}

Short-circuit evaluation (SCE) refers to the omission of evaluation of operand(s) in a boolean
expression when the result of the boolean expression would not be changed by their evaluation(s).
SCE is supported by C-family (also the mainstream) languages such as C, C++, \java, \python and
many others.
SCE is an interesting feature for better or for worse.
Unprepared developers may be caught by surprise.
Inadvertent usage may results in unexpected results.
We will use a couple of examples to illustrate implementation skills around SCE\@.

Consider the following coding requirements:
\begin{quote}
    ``Implement a method, \code{next}, to fetch the next item from a source
    which is irrelevant to this problem and
    a method \code{procure} has been made available to fetch an element from the source and
    stored it in a member variable \code{Optional<Object> cache} as
    \lstinputlisting{if-else/getNext_procure.java}
    The cache shall be cleared with each successful invocation of method \code{next}.
    ''
\end{quote}
This is a subtask of an implementation of a class.
Presumably, the algorithm to solve this problem would be
\begin{quote}
    \begin{enumerate}
        \item Check \code{cache} to see if it is empty, if so
        \begin{enumerate}
            \item Invoke \code{procure} to attempt to populate \code{cache}
            \item If attempt of procuring an item is unsuccessful, throw ``No element is left''
            exception
        \end{enumerate}
        \item Get the item from cache and clear the cache
        \item Return the item
    \end{enumerate}
\end{quote}

Let's see how we would implement the algorithm in \java in the absence of SCE mechanism.
\lstinputlisting[caption={Implementation of \code{next} in the absence of SCE mechanism},label={lst:next-without-sce}]
{if-else/getNext.java}
It seems pretty simple code and is not bad.
But as we all know \java supports SCE and with it one can readily implement
the following logic instead
\lstinputlisting[caption={Implementation of \code{next} method taking advantage of SCE in\java.},
label={lst:next-sce}]
{if-else/getNext_shortCircuit.java}
Here is how SCE works in ~\autoref{lst:next-sce}.
When \code{cache} contains an item, the second part in the
\code{if}-expression is \emph{ignored}, namely, the method \code{procure} is \emph{not} invoked.
The entire \code{if}-statement is skipped and the item in \code{cache} is all the business.
When \code{cache} is empty, the \code{if}-expression continues with the evaluation of the
second part, that is, \code{procure} is invoked.
Depending on the return value of the invocation, exception is thrown if \code{false} or control
is passed to subsequent statement after \code{if} statement.

Comparing ~\autoref{lst:next-sce} with ~\autoref{lst:next-without-sce}, we saw an elegant,
succinct implementation of otherwise messy code.

What can we learn about short-circuit mechanism from this example?
\begin{itemize}
    \item Order of operands matters, \textit{e.i.}, \code{a && b} is not equivalent to
    \code{b && a}.
    \item
\end{itemize}

How to consciously seek opportunities of SCE mechanism to simplify one's code?
Firstly awareness will put you in an advantageous position of mastering elegant coding skills.
You need to know the ins and outs of SCE.
Next one needs to consciously modularize one's design, creating reusable, meaningful
single-responsible subroutines to be invoked at higher level.
We will talk more about function designs in ~\fullref{ch:reusableFunctions}
For example one need to create functions \code{procure} with a well-defined return value
indicating if procurement of an item is successful or not for use in the short-circuit boolean
expression in ~\autoref{lst:next-sce}.
One need to be aware of other characteristics such as idempotence and mutability of the
function invocation.

%Before ending this section, let us look at a more subtle case code snippet
%~\autoref{lst:ternary-boolean-expr}.
%Assuming that \code{boolArr} is an array of \code{boolean} type,
%~\autoref{lst:ternary-boolean-expr}
%% side-by-side two columns
%\newline\begin{minipage}[t]{.45\textwidth}
%             \lstinputlisting[caption={Ternary boolean expression},
%             label={lst:ternary-boolean-expr},
%             firstline=1,lastline=1]
%             {if-else/equivalent_statements_ternary_expr.java}
%\end{minipage}\hfill
%\noindent\begin{minipage}[t]{.45\textwidth}
%             \lstinputlisting[caption={Simplified boolean expression},
%             label={lst:simplified-boolean-expr},
%             firstline=3,lastline=3]
%             {if-else/equivalent_statements_ternary_expr.java}
%\end{minipage}


\section{Case-stack}\label{sec:case-stack}

As a programming language feature,
it is not unfamiliar to many a sequence of consecutive conditional statements
starting with \code{if} clause, followed by zero, one, or more of
\code{else if} blocks (\code{elif} for \python), and optionally with
\code{else} block, \textit{e.g.}, ~\autoref{lst:simpleBranching}.
For the sake of argument, we refer to such construct as `case-stack'.
Case-stack is an important pattern to implement a wide range of algorithms.
In combination with loop constructs or within recursive functions,
Case-stack forms an effective pattern for implementing complicated logic.
Its applications are throughout this book.
Please refer to examples in ~\fullref{ch:loops}, ~\fullref{ch:recursion}, and the chapters in
~\fullref{prt:case-studies} for more discussion.

Unfortunately, little is mentioned about this patterns in programming references.
Beginning and intermediate developers often stumble upon this without conscious choice, lest to
master the best practice of it.
In this section, so we hope to give it a fair coverage over some of the important aspects of a
case-stack and the best practices around it.

When faced with multi-branch switching,
a common way of implementation uses nested \code{if..else} statements such as the example given at
the beginning of this chapter ~\autoref{lst:nestedBranching}.
As switching fanout grows, the nesting level goes deeper, indentation grows messier, and
the logic quickly becomes unwieldy.
As an alternative for better readability, one can organize multi-branch switching using a
case-stack as ~\autoref{lst:simpleBranching}.
For some special cases, \code{switch}-statement may be used equally well, if supported by the
programming language.
For simplicity, we limit our discussion to \code{if..else}-statement here.

The execution model of a case-stack is like walking up a staircase---the program
successively evaluates the conditional expression one at a time and builds up the cases.
The conditional expression of the leading \code{if}-block is the first step of check.
If it is evaluated to \code{true}, then the \code{if}-block is executed
and the rest of the stack is skipped.
Otherwise, the next conditional expression is evaluated
and so on until the end of the case-stack.
Note that in the presence of \code{return} and \code{break},
the case-stack may be simplified with the removal of the \code{else} keyword.
For examples, take a look at ~\autoref{lst:if-stack-equals}, ~\autoref{lst:if-stack-lca},
and ~\fullref{lst:matrix-spiral-term} in ~\fullref{ch:recursion}.

First property of a case-stack is the strict order of conditional statements which cannot be
altered once completed.
The order of conditional statements takes careful design for succinct code,
optimal performance or both.
The second property is the cumulative effect of conditional expressions.
We are going to talk these properties below.

\begin{table}[tbh]
    \begin{center}
        \caption{\label{tab:cumulative-tests-case-stack}Cumulative effect of tests in a case
        stack.}
        \begin{tabular}{p{7cm}|p{7cm}}
            Case stack & Equivalent conditional statements \\
            \lstinputlisting{if-else/case_stack_example.java}
            &
            \lstinputlisting{if-else/case_stack_taken_apart.java} \\
        \end{tabular}
    \end{center}
\end{table}
When taken apart, the effect of tests of a case-stack
is cumulative---a later case combines all previous conditional expressions up to its own.
This is illustrated by the example of ~\autoref{tab:cumulative-tests-case-stack}.
Note that unlike the strict ordering of the conditional statements in the left column,
the three \code{if}-blocks in the right column are not logically coupled and
thus may be swapped assuming that the execution of \code{statement1}, \code{statement2},
\code{statement3}, do not affect one another.

When debugging a case-stack for logic fallacies and unfounded assumptions, for example,
one needs to keep this cumulative effect in mind: all the assumptions about the
the current line of code should have been tested by an earlier conditional expression.
A missing test will likely result in exception in subsequent execution.
From this perspective, the right column in ~\autoref{tab:cumulative-tests-case-stack} would be a
good aid for debugging---due to its explicit conditional expressions, it is more
straightforward to spot untested assumptions made for a later case than its case-stack
counterpart.

\begin{figure}[htb]
    \lstinputlisting[label={lst:if-stack-equals},
    caption={Example implementation overriding \code{equals} in a subclass}]
    {if-else/equals.java}
\end{figure}
As a real-world example to demonstrate the effectiveness of case-stack in handling
multi-branching logic,
let's look at a typical \code{equals} method in a subclass that overrides
the default implementation in \java.
As shown in ~\autoref{lst:if-stack-equals},
The first test we do is the nullity of the argument.
Only if it is not null do we move on to the second test---if it's of type \code{MyClass}.
Only if it is not null and is of type \code{MyClass} do we cast it to type \code{MyClass}
and perform pairwise equality tests between member variables of
\code{this} instance and \code{that} instance.
Once again, you see how the program progresses up a `staircases' step-by-step.

\begin{figure}[htb]
    \lstinputlisting[label={lst:if-stack-lca},
    caption={An implementation of finding lowest common ancestor in a binary search tree.}]
    {if-else/lca.java}
\end{figure}
Earlier we pointed out the strict ordering property of case-stack.
To illustrate how the order of tests affects the construct of case-stack,
let's consider the problem:
\begin{quote}
    ``determining the lowest common ancestor of two values in a binary search tree.''
\end{quote}
By the nature of binary search tree and the problem requirements, the number of `special' cases
is large.
If not designed properly, the code of handling all the cases can blow up like a nightmare.
However, with a properly designed case-stack as we are going to show next, a couple of lines of
code can handle the myriad possibilities like a breeze.
Take a look at ~\autoref{lst:if-stack-lca}.
The cases are built from ground up like this:
\begin{enumerate}
    \item Test nullity;
    \item Test for left branching;
    \item Test for right branching;
    \item Test for existence of keys;
    \item Final case;
\end{enumerate}
Note later statements are built on the assertions of previous tests---\textit{e.g.}, test for
left/right branching requires the non-null attributes of root otherwise the code will throw
\code{NullPointerException}.
Also note that as the cases are being built up, the solution space is also being halved
and ruled out much like that of a binary search.
Some source states that when doing tests to rule out cases in settings such as a case-stack, start
with the simplest one.
This observation is quite superficial.
The real criterion is to do the test that contains most entropy, the test that is most
informative, the test that is hinged with most stake first.
That way, the maximal solution space can be chopped off at the earliest possible time.
Generally speaking, building up the cases of an case-stack is ultimately empirical.
Even the experienced developers may have to make a couple of attempts to get the optimal design.
So go ahead and try this out this in your own practice and sharpen your instinct about it.

\section{Conditional statements in a loop}\label{sec:branchingInLoop}

Conditional statements in a loop is one of the most important application.
With the keywords \code{break} and \code{continue}, conditional statements
have more moving parts, therefore logic simplification in a loop is even more challenging.
On the other hand, this also means skilled developers have new weapons to wield when aiming at
simple, succinct logic.

Practical programming tips for loops will be more fully discussed ~\autoref{ch:loops}.
Many of the topics overlap between here and there,
so we are going to be brief here and with emphasis on conditional constructs.
For fuller discussions, we direct readers there.

For a really good example of conditional application in a loop,
the ~\fullref{lst:silhouette-skyline} presented in ~\fullref{ch:loop-thinning}
would make an excellent example.

Like a complex protein molecule, the conditional statements are assembled into a case-stack in
resemblance of hierarchical structures and together they perform the processing in a stateful loop.
